\section{Constant cochains, modification, and stable transport}
\label{sec:constants}

We compare constants on the division and matrix sides before choosing their
generators.  This distinction matters: an abstract isomorphism of exterior
algebras would not identify the cup action in
Proposition~\ref{prop:full-frame-comparison}.  The comparison below is induced
by pullback from a single modification stack.  It also applies to finite
Postnikov coefficients, and hence transports continuous sphere-cochain
representatives.

Throughout this section, $1\leq a<p-1$, and
\[
 K_a=\GL_a(\Zp),\qquad
 I_a=\{g\in K_a:g\bmod p\text{ is upper triangular}\},\qquad
 P_a=\cO_{D_a}^{\times},
\]
where $D_a/\Qp$ has invariant $1/a$.  All group cohomology is continuous.
Coefficients such as $\Zp$ carry their profinite topology; finite coefficients
carry the discrete topology.  We use $C^*_{\mathrm{cts}}(G;E)$ for the spectrum
of continuous cochains with constant spectrum $E$, and
$\RGamma_{\mathrm{cts}}(G;M)$ for its complex-valued counterpart.  Finite
Postnikov spectra here have finitely many nonzero homotopy groups, each a
finite $p$-group.  Cochains with these coefficients are formed first; completed
coefficients will be obtained by the explicit inverse limits below.

\subsection{Integral and finite-coefficient cohomology}

\begin{theorem}[Constant cohomology]
\label{thm:constant-cohomology}
For $G=K_a,I_a$, or $P_a$, the groups $H^q(G,\Zp)$ are finite free
$\Zp$-modules, vanish for $q>a^2$, and have Poincar\'e polynomial
\begin{equation}
 \sum_q\operatorname{rank}_{\Zp}H^q(G,\Zp)z^q
       =\prod_{i=1}^a(1+z^{2i-1}).
 \label{eq:const-poincare}
\end{equation}
There are exterior-algebra presentations over $\Zp$, and over every finite
extension of $\Fp$, with generators in degrees $2i-1$, $1\leq i\leq a$.
For every $r\geq1$, reduction induces an isomorphism
\begin{equation}
 H^q(G,\Zp)/p^r\ \xrightarrow{\ \sim\ }\ H^q(G,\Z/p^r).
 \label{eq:const-reduction}
\end{equation}
The actual restriction maps
\[
 H^*(K_a,\Zp)\longrightarrow H^*(I_a,\Zp),\qquad
 H^*(K_a,\Z/p^r)\longrightarrow H^*(I_a,\Z/p^r)
\]
are isomorphisms.  These assertions do not select the generators: the
compatible stable generators will be constructed in
Theorem~\ref{thm:stable-generators}.
\end{theorem}

\begin{proof}
We specify the finite-coefficient input.  Theorem~1.1 of
\cite{DLH}, comprising Propositions~5.5 and~5.7 there, computes the cohomology
of an unramified split Chevalley group and proves restriction to its Iwahori
to be an isomorphism when $p>h+1$, where $h$ is the Coxeter number.
For $\GL_a/\Qp$, $a\geq2$, this is $h=a$; the degree of the ground field and
its ramification index are both one.  Thus its hypothesis is exactly
$a<p-1$, and its exterior degrees are $1,3,\ldots,2a-1$.
For the division algebra, \cite[\S5.2, Proposition~5.12]{DLH} applies to an
unramified ground field and a division algebra of invariant $1/a$ with
$a<p-1$.  It gives the same exterior algebra over the residue field of
$D_a$, which here is $\F_{p^a}$.  Extension of the finite constant field
commutes with continuous cochains.  In particular, that calculation gives
the stated degreewise dimensions over $\Fp$.  For $a=1$ all three groups
are $\Zp^{\times}=\mu_{p-1}\times(1+p\Zp)$, so the assertion follows
directly from the two-term resolution for $\Zp$ and exactness of
$\mu_{p-1}$-invariants.

We include the integral deduction.  None of these groups has an element of
order $p$.  Such an element in $\GL_a(\Qp)$ would require the cyclotomic
polynomial $\Phi_p$, of degree $p-1$, in its rational representation; in
$D_a^{\times}$ it would embed $\Qp(\zeta_p)$ as a subfield of a division
algebra of degree $a$.  Both are impossible for $a<p-1$.  The compact
$p$-adic analytic group theorem therefore gives $p$-cohomological dimension
$a^2$.  Its completed group algebra is Noetherian, and the trivial module
has a bounded resolution by finitely generated projectives.  One can obtain
finite generation of the resolution terms from Noetherianity and truncate
at the cohomological dimension.  The compact-module form of these
analytic-group results is recalled in
\cite[Sections~1.1--1.2]{Venjakob}; see also \cite{Lazard,SerreGalois}.
In particular, $H^q(G,\Zp)$ is finitely generated over $\Zp$, and rational
base change computes continuous rational cohomology.

Write $r_q$ for its rank and $t_q$ for the number of cyclic summands in its
finite $p$-primary torsion subgroup.  The coefficient sequence gives
\[
 0\longrightarrow H^q(G,\Zp)/p
 \longrightarrow H^q(G,\Fp)
 \longrightarrow H^{q+1}(G,\Zp)[p]\longrightarrow0,
\]
and consequently
\begin{equation}
 \dim_{\Fp}H^q(G,\Fp)=r_q+t_q+t_{q+1}.
 \label{eq:const-uct}
\end{equation}
Rational Lazard comparison identifies the rational cohomology with the
invariant Lie-algebra calculation.  The matrix Lie algebra is
$\mathfrak{gl}_a$; after a splitting field extension the division Lie
algebra is the same Lie algebra.  Its cohomology is exterior in degrees
$1,3,\ldots,2a-1$, and conjugation acts trivially
\cite[Proposition~3.8.1 and Lemma~3.8.2]{BSSW}.  The rational and
mod-$p$ dimensions therefore agree in every degree.  Equation
\eqref{eq:const-uct} forces $t_q=t_{q+1}=0$ for every $q$.
The coefficient sequence for $p^r$ now proves
\eqref{eq:const-reduction}.  Restriction to $I_a$ is an isomorphism
integrally as well, since its reduction modulo $p$ is an isomorphism
between finite free modules of equal rank.

For completeness, the finite-field extension in the division calculation
does not conceal an algebra-descent assertion.  Choose over $\Fp$ a
homogeneous basis of the indecomposable quotient of the positive-degree
cohomology algebra.  After extension to $\F_{p^a}$ there is precisely one
basis element in each of the indicated odd degrees.  Their lifts generate
the original graded algebra by induction on degree; their squares vanish
because $p$ is odd.  The resulting surjection from the exterior algebra is
an isomorphism by the degreewise dimension calculation.  Lift these
generators to integral cohomology using \eqref{eq:const-reduction}.
Graded commutativity and torsionfreeness give a map from the integral
exterior algebra, whose reduction modulo $p$ is an isomorphism.
Nakayama's lemma in each degree proves the integral presentation.
\end{proof}

\begin{remark}
\label{rem:const-galois}
The coefficient-field automorphisms used to pass from the ordinary to the
extended stabilizer act trivially on $H^*(P_a,\Zp)$.  Indeed, on the
division Lie algebra they become inner automorphisms after splitting, so
they act trivially on rational cohomology.  Torsionfreeness then makes
their integral action trivial.  This concerns the constant cohomology
classes; equivariant sphere representatives can subsequently be obtained
by transfer across a finite extension of degree prime to $p$.
\end{remark}

\subsection{The subgroup of determinant valuation zero}

Set
\[
 J_a=\{g\in\GL_a(\Qp):v_p(\det g)=0\}.
\]
This condition retains a lattice in the determinant of a rational frame,
without requiring a lattice in the whole frame.

\begin{lemma}[Building restriction]
\label{lem:const-building}
Restriction along $K_a\subset J_a$ induces an equivalence
\[
 C^*_{\mathrm{cts}}(J_a;E)\ \xrightarrow{\ \sim\ }
 C^*_{\mathrm{cts}}(K_a;E)
\]
for every finite Postnikov $p$-primary constant coefficient spectrum $E$.
It also induces isomorphisms with constant coefficients $\Zp$ and
$\Z/p^r$.
\end{lemma}

\begin{proof}
For $a=1$ the groups are equal.  Otherwise let $\mathcal B_a$ be the
reduced Bruhat--Tits building.  The group $J_a$ preserves vertex types,
acts without inversions, and has a closed chamber as fundamental domain.
The stabilizer $K_v$ of a vertex is a conjugate of $K_a$.  Indeed a matrix
stabilizing the homothety class of a lattice is a scalar times a lattice
automorphism, and determinant valuation zero forces the scalar valuation
to be zero.  Every face stabilizer $K_\sigma$ is compact and lies between
the Iwahori of the chamber and one of its vertex stabilizers.

The index $[K_\sigma:I_a]$ is prime to $p$: after reduction it is a flag
count for the corresponding Levi blocks, and each such flag count is
congruent to one modulo $p$.  Transfer therefore makes
$H^*(K_\sigma,\Fp)\to H^*(I_a,\Fp)$ injective.  The restriction
$H^*(K_v,\Fp)\to H^*(I_a,\Fp)$ is an isomorphism by
Theorem~\ref{thm:constant-cohomology}, and factors through this injection.
It follows that the face restriction is an isomorphism too.  Identifying
every face group with $H^*(I_a,\Fp)$ by its actual restriction makes all
incidence maps identities.

The equivariant cellular resolution of the contractible building gives
\begin{equation}
 E_1^{r,s}=\bigoplus_{\substack{\sigma\text{ in the chamber}\\
                              \dim\sigma=r}}
                 H^s(K_\sigma,\Fp)
       \ \Longrightarrow\ H^{r+s}(J_a,\Fp).
 \label{eq:const-building-ss}
\end{equation}
The sum is finite, the chamber has dimension $a-1$, and the compact face
groups have cohomological dimension $a^2$.  Thus this is a bounded
spectral sequence.  Its rows are the ordinary cochain complex of a
simplex with constant coefficient $H^s(I_a,\Fp)$.  They have only degree
zero cohomology.  The resulting edge map is restriction to a vertex;
this proves the assertion for $H\Fp$.

A finite abelian $p$-group has a finite filtration with quotients $\Fp$.
Coefficient exact sequences and then Postnikov fibre sequences prove the
assertion for every stated $E$.  This proves a comparison of the actual
restriction maps, natural in $E$.  The integral assertion follows by
taking the inverse limit of the finite-coefficient comparisons.  The
groups on the compact side are finite at each finite coefficient level,
so the cohomological inverse systems are Mittag--Leffler.  The same is
then true on the $J_a$ side.
\end{proof}

\subsection{The common modification stack and its arithmetic}

Fix $n<p-1$.  Choose a finite unramified extension $K/\Qp$ over which the
Honda endomorphisms in use are defined, and a completed algebraic closure
$C$ of $K$.  All spaces in the next display are geometric diamonds over
$C$; we retain their $\Gal(\overline K/K)$-descent.  Let
\[
 \Omega^{n-1}_C=\mathbb P^{n-1}_C\setminus
       \bigcup_H H_C,
\]
where $H$ ranges through rational $\Qp$-hyperplanes.

\begin{lemma}[Modification presentations]
\label{lem:const-modification-stack}
There is a common modification stack, with its determinant lattice,
\begin{equation}
 \mathcal M_n\simeq[\mathbb P^{n-1}_C/P_n]
             \simeq[\Omega^{n-1}_C/J_n].
 \label{eq:const-modification-stack}
\end{equation}
The two presentations commute with the arithmetic action over $K$.
In particular, arithmetic acts trivially on the continuous group
cochains of $P_n$ and $J_n$.  Their structural maps to geometric
cochains, constructed in the next subsection, are
arithmetic-equivariant.
\end{lemma}

\begin{proof}
Use the convention that $V_n$ has slope $1/n$.  A point of projective
space specifies a length-one quotient at the untilt divisor,
\[
 0\longrightarrow W\longrightarrow V_n'
          \longrightarrow i_*\mathcal L\longrightarrow0,
\]
where $V_n'$ is a form of $V_n$.  The kernel has rank $n$ and degree
zero.  If a proper subbundle of rank $r$ had positive degree, its
inclusion in the stable bundle $V_n'$ would give degree strictly less
than $r/n<1$, a contradiction.  Thus $W$ is semistable of slope zero.
The relative slope-zero equivalence identifies it with a rank-$n$
$\Qp$-local system.

Conversely, an upper modification of a trivial rank-$n$ bundle has
nonnegative slopes and total degree one.  Nonnegativity follows because
a negative-slope quotient would receive no map from the trivial bundle,
whereas that bundle has only a torsion cokernel.  The positive part is
therefore a single stable bundle of degree one; the rest, if present,
has slope zero.  A slope-zero quotient is precisely a rational linear
functional annihilating the modification direction.  Its absence is
exactly the condition that the direction lie in $\Omega^{n-1}_C$.
On that locus the middle bundle is of type $V_n$.

The relative classification of these bundles and the slope-zero/local-system
equivalence make the frames torsors in the $v$-topology, so the two
descriptions are inverse on families as well as on geometric points.
This is the basic EL modification description underlying the duality
of \cite[Theorem~7.2.3]{ScholzeWeinstein}; the slope argument above
specifies its basic locus in the present case.

We explain the determinant reduction.  Fix a lattice in the rank-one
slope-zero local system
$\det(V_n)(-\infty)$.  Require the determinant of the frame of $W$ to
carry $\Zp$ to that lattice.  Its permitted rational frame changes are
exactly $J_n$.  On the positive-bundle side the corresponding condition
is $v_p(\operatorname{Nrd}(g))=0$ on $D_n^{\times}$, whose subgroup is
$P_n$.  Hence the same determinant condition cuts out the two groups
in \eqref{eq:const-modification-stack}.  The determinant-matching
construction is the one appearing, for the positive-height extension
towers, in \cite[Theorem~2.0.1 and Theorem~2.6.3]{BMR}.

Choose the positive Honda isocrystal over the unramified field containing
its endomorphism field.  Arithmetic then acts on the period factor and
commutes with its division-algebra endomorphisms.  Subtracting the untilt
divisor in its determinant produces a rank-one $\Qp$-local system with
compact arithmetic image; equivalently, its integral Tate lattice is
preserved.  More intrinsically, the valuation of an arithmetic
determinant action is a continuous homomorphism from a profinite group
to $\Z$ and hence is zero.  The chosen determinant component can
therefore be retained arithmetically.  On the split-frame directions
this is the usual arithmetic action on Drinfeld space: changing the
trivialization of the determinant only changes a line scale, which
disappears in projective directions.  Thus both torsor presentations
and both structural maps commute with arithmetic.
\end{proof}

\subsection{Finite constant sections and geometric descent}

The two presentations of $\mathcal M_n$ now specify the comparison to
be made.  We need its structural maps to preserve the actual constant
classes and their products, and we need arithmetic to separate these
classes from higher geometric cohomology.  We first construct the
maps on finite coefficients.

The compact-group descent calculation is
Lemma~\ref{lem:solid-rows}.  Its geometric parameter construction and
completed-bar comparison apply to both period spaces above.  For
$J_n$, use the finite chamber resolution of
Lemma~\ref{lem:const-building} and apply that calculation to each
compact face stabilizer.  There are finitely many faces, so this adds
a finite filtration.  In particular, a commuting arithmetic scalar
on a geometric row stays the same scalar on all its group-cohomology
rows.  This retains the profinite parameters of the Steinberg duals
that will occur for Drinfeld space.

We make explicit the map from the separately defined group cochains to
geometric cochains.  For a small $v$-stack $X$, write
$C_v^*(X;E)=R\Gamma(X_v,\underline E)$ for the cochains of the constant
sheaf associated to a finite Postnikov $p$-primary spectrum $E$.  For a
compact profinite set $T$, put
\[
 \operatorname{LC}(T;E)
   =\operatorname*{colim}_{T\twoheadrightarrow T_i}
       \prod_{t\in T_i}E,
\]
where $T_i$ runs over the finite clopen quotients.  For the locally
profinite group powers used here, take the product of these objects over
a disjoint compact-open decomposition.  Each compact piece meets only
finitely many pieces of another such decomposition, so common refinement
identifies the result.  For abelian coefficients the same notation means
the locally constant function module, degreewise for complexes.  On each
finite clopen partition the
constant-section maps for $E$ on the inverse-image pieces define
\[
 \operatorname{LC}(T;E)\longrightarrow
 C_v^*(X\times\underline T;E).
\]
They commute with refinement.  If a locally profinite group $G$ acts on
$X$ over a small $v$-stack base $\mathcal B$, the action nerve therefore
gives, for the trivial action on $E$,
\[
 \begin{split}
 \eta_{X,G,E}:\ C^*_{\mathrm{cts}}(G;E)
   &=\operatorname{Tot}_{a\geq0}\operatorname{LC}(G^a;E)\\
   &\longrightarrow
     \operatorname{Tot}_{a\geq0}
       C_v^*(X\times\underline{G^a};E)
       \simeq C_v^*([X/G];E).
 \end{split}
\]
The last equivalence is \v{C}ech descent for the quotient.  This
constant-section construction is valid over every $\mathcal B$.
If $\phi:H\to G$ is continuous and $f:X'\to X$ is $\phi$-equivariant
over a base map, then the induced quotient map satisfies
\[
 [f]^*\eta_{X,G,E}
   =\eta_{X',H,E}\operatorname{res}_{\phi}.
\]
The equality holds on the nerve terms and commutes with all faces and
degeneracies.  The same termwise construction commutes with coefficient
maps and, for any specified pairing $E\wedge F\to E'$, with the
corresponding cup maps.  It is consequently natural under arithmetic
base automorphisms as well.

Write $B_{\mathcal B}G=[\mathcal B/\underline G]$ for the relative
classifying stack.  Its universal map $\eta_{\mathcal B,G,E}$ pulls back
to $\eta_{X,G,E}$ along $[X/G]\to B_{\mathcal B}G$.  If
$\mathcal T\to Z$ is a $G$-torsor with classifying map $\beta$, its
\v{C}ech nerve is $\mathcal T\times\underline{G^a}$, so under
$[\mathcal T/G]\simeq Z$ the map $\beta^*\eta_{\mathcal B,G,E}$ is
$\eta_{\mathcal T,G,E}$.  This identity commutes with base change of the
torsor.  For an $H$-torsor $\mathcal T$ and a homomorphism $\phi:H\to G$,
it also gives, on the common quotient $Z$,
\[
 \eta_{\mathcal T\times^H G,G,E}
   =\eta_{\mathcal T,H,E}\operatorname{res}_{\phi}.
\]
In geometric occurrences of $BG$ in this paper, the base is the one in
the diagram at issue, and a classifying pullback of a continuous class
$c$ means the pullback of its universal image
$\eta_{\mathcal B,G,E}(c)$.  The structural constant actions on supported
coefficients are the cup actions of these same finite unit images on the
common \v{C}ech nerves used for support.

This definition agrees with the geometric parameter cochains used in
Lemma~\ref{lem:solid-rows}.
For locally spatial $Y\times\underline T$ with $T$ profinite,
\cite[Proposition~14.10]{ScholzeDiamonds} identifies finite constant
\'{e}tale cochains with their $v$-cochains.  For a finite abelian
$p$-group $M$, use the coefficient ring $\Z/p^r$ annihilating $M$ in
that proposition; its coefficient ring is arbitrary.  Finite Postnikov
induction extends the identification from $HM$ to $E$.  Write
$\mathscr C_{Y,\Fp}$ for the derived solid parameter object associated
to $Y$ in that construction.  The constant sections give a unit
$\underline{\Fp}\to\mathscr C_{Y,\Fp}$, and the just specified
\'{e}tale--$v$ comparison identifies the map on compact bar invariants
induced by this unit with $\eta_{Y,G,H\Fp}$.  Thus the unit is fixed
before the comparison of the bar resolution with $Q_\bullet$.

For a geometric point $B_C=\operatorname{Spd}C$, with $C$ complete and
algebraically closed, the universal map is an equivalence when $G$ is
compact.  Indeed, the example after Definition~7.8 in
\cite{ScholzeDiamonds} expresses $B_C\times\underline T$, for profinite
$T=\varprojlim T_i$, as the limit of the finite disjoint unions
$B_C\times T_i$.  Proposition~7.16 there makes the $C$ point strictly
totally disconnected, so its \'{e}tale covers split and its finite
constant cohomology is concentrated in degree zero.  Propositions~14.9
and~14.10 then give the actual constant-section equivalence
\[
 \operatorname{LC}(T;HM)
    \simeq C_v^*(B_C\times\underline T;HM)
 \qquad(M\text{ a finite abelian }p\text{-group}).
\]
Finite Postnikov induction gives the same assertion for $E$.  The
equivalences are natural in $T$ and in automorphisms of $C$; taking them
on the compact $G$-nerve proves the claim about
$\eta_{B_C,G,E}$.  We use this compact comparison for the literal
relative classifying interpretation of $P_n$ and $K_n$.  For $J_n$ and
the upper block group below, the universal images together with the
finite chamber and inflation calculations provide the required maps.
For the modification stack $\mathcal M_n$ above, the unadorned
$C^*(\mathcal M_n;E)$ denotes $C_v^*(\mathcal M_n;E)$ and
$H^q(\mathcal M_n,\Fp)=\pi_{-q}C_v^*(\mathcal M_n;H\Fp)$.
In what follows, constant cochains on $BP_n$ and $BJ_n$ mean these
continuous sources with their universal images.

\subsection{Arithmetic separation and continuous stable transport}

Choose $\gamma\in\Gal(\overline K/K)$ such that
$\bar\chi_p(\gamma)=u$ generates $\Fp^{\times}$.  Such a choice is
possible because $K/\Qp$ is unramified.  For a vector space with an
endomorphism annihilated by a polynomial whose roots lie in
$\Fp^{\times}$, its generalized weight-zero part means its
$(\gamma-1)$-primary summand.  Nilpotent extensions are included in this
terminology.

\begin{proposition}[Transport of constants]
\label{prop:modification-transport}
The structural maps from \eqref{eq:const-modification-stack} identify
\begin{equation}
 H^*(P_n,\Fp)\ \xrightarrow{\ \sim\ }
 H^*(\mathcal M_n,\Fp)_{(1)}
 \xleftarrow{\ \sim\ }H^*(J_n,\Fp),
 \label{eq:const-weight-zero}
\end{equation}
where $(1)$ denotes generalized eigenvalue one for $\gamma$.
These are identifications by the actual pullback maps and preserve cup
products.  Together with Lemma~\ref{lem:const-building}, they induce
natural equivalences, for every finite Postnikov $p$-primary spectrum $E$,
\begin{equation}
 C^*_{\mathrm{cts}}(P_n;E)\ \simeq\
 C^*(\mathcal M_n;E)_{(1)}\ \simeq\
 C^*_{\mathrm{cts}}(J_n;E)\ \simeq\
 C^*_{\mathrm{cts}}(K_n;E).
 \label{eq:const-postnikov-transport}
\end{equation}
The equivalences commute with coefficient maps, products, and Hurewicz
maps.  They extend to the $p$-complete sphere by the Postnikov and
coefficient limits defining continuous cochains.

There is also the following compatibility on the independent extension
locus $Z=(N_t^m)_{\mathrm{ind}}$, $m=n-t$.  After fixing the connected
quotient frame, let
\[
 b:[Z/K_m]\longrightarrow BP_n,\qquad
 q:[Z/K_m]\longrightarrow BK_m
\]
be the positive middle-bundle classifying map and the structural map.
If $c_P$ and $c_K$ correspond under
\eqref{eq:const-postnikov-transport}, then
\begin{equation}
 b^*c_P=q^*\operatorname{res}_{K_m}^{K_n}(c_K),
 \qquad K_m\longrightarrow K_n,\quad
             g\longmapsto\operatorname{diag}(g,1_t).
 \label{eq:const-block-action}
\end{equation}
This equality holds for finite Postnikov coefficients, for their stable
limits, and for the resulting ordinary cup actions.  In particular the
constant action in Proposition~\ref{prop:full-frame-comparison} is the
actual block-restriction action.
\end{proposition}

\begin{proof}
\emph{Arithmetic separation and the specified degree-zero edges.}
We first prove \eqref{eq:const-weight-zero}.  The projective-space
calculation is
\[
 H^{2j}(\mathbb P^{n-1}_C,\Fp)=\Fp(-j),\quad 0\leq j<n,
 \qquad H^{2j+1}(\mathbb P^{n-1}_C,\Fp)=0.
\]
Its geometric automorphisms act trivially on these lines, and $\gamma$
acts by $u^{-j}$.  For Drinfeld space, apply the integral and
finite-coefficient theorem \cite[Theorems~1.1 and~5.1]{CDN} with ground
field $\Qp$, matching the rational $\Qp$-hyperplanes defining
$\Omega_C^{n-1}$, and then restrict its arithmetic action to
$\Gal(\overline K/K)$.  It gives compatible topological isomorphisms
\begin{equation}
 \begin{split}
 H^j_{\mathrm{\acute et}}(\Omega^{n-1}_C,\Zp(j))
       &\simeq\operatorname{Sp}_j(\Zp)^*,\\
 H^j_{\mathrm{\acute et}}(\Omega^{n-1}_C,\Fp(j))
       &\simeq\operatorname{Sp}_j(\Fp)^*,
                  \qquad 0\leq j<n.
 \end{split}
 \label{eq:const-cdn}
\end{equation}
Here $\operatorname{Sp}_j$ is the integral generalized Steinberg
representation, the star is its continuous integral or finite-field
dual, and arithmetic acts trivially on that factor.  Cohomology vanishes
above $n-1$.  Thus the untwisted $j$th group has exactly the arithmetic
scalar $u^{-j}$; degree zero is precisely the constants.  In particular,
\eqref{eq:const-cdn} is an integral and finite-coefficient statement.
We neither reduce a rational pro-\'{e}tale calculation nor discard
analytic terms by tensoring it with $\Fp$.

Apply geometric descent to the two presentations of $\mathcal M_n$.
Lemma~\ref{lem:solid-rows} shows that every group-cohomology row arising
from geometric degree $j$ has the same arithmetic scalar $u^{-j}$.
For the projective presentation its geometric degree is $2j$; for
the Drinfeld presentation it is $j$.  Both filtrations are bounded,
using the compact group resolution on the first side and the finite
chamber together with compact face resolutions on the second.  Since
\[
  1\leq j\leq n-1<p-1,
\]
all the higher geometric rows have eigenvalue different from one.
An equivariant differential between generalized summands with relatively
prime eigenvalue polynomials is zero.  Exactness of primary decomposition
therefore shows that the generalized-one part of the abutment is exactly
the geometric degree-zero row.

We identify its edge using the unit
$\underline{\Fp}\to\mathscr C_{Y,\Fp}$.  For each of
$Y=\mathbb P_C^{n-1}$ and $Y=\Omega_C^{n-1}$, its induced point map
$\Fp\to H^0_{\mathrm{\acute et}}(Y,\Fp)$ is an isomorphism: the
displayed row calculations make the target one-dimensional, and the
constant section $1$ is nonzero.  For each compact group in the
descent, apply the finite projective resolution $Q_\bullet$ of
Lemma~\ref{lem:solid-rows}.  Each evaluated internal-Hom term is an
idempotent summand of finitely many copies of the point value, as in
\eqref{eq:const-solid-resolution}.  The unit commutes with these
idempotents and the differentials, so it is an isomorphism on the
bottom-row computing complexes.  The bar comparison identifies the
map so obtained with $\eta_{Y,G,H\Fp}$ from continuous cochains.  On
the Drinfeld side make the same check on each compact face and then
take the finite chamber: naturality of the unit with respect to face
restrictions holds before the projective resolutions are chosen.
Thus both degree-zero edges are the specified finite constant-section
maps.  This proves \eqref{eq:const-weight-zero}, with its specified
maps.  The generalized-one part is closed under products, and these
unit maps are multiplicative, proving the assertion about cup products.

\emph{Finite Postnikov coefficients and the sphere limit.}
The finite-coefficient comparison must now retain actual lifts.
Put
\[
 Q(T)=\prod_{j=1}^{n-1}(T-\widetilde{u^{-j}})\in\Z[T],
\]
where the tildes are integer lifts; the empty product at $n=1$ is one.
Then $Q(1)$ is a $p$-adic unit.  The preceding bounded filtrations imply
that, on the finite-coefficient cohomology, some powers of $T-1$ and
$Q(T)$ give the two coprime primary factors for $T=\gamma$.
Coefficient extensions and Postnikov fibre sequences retain this
property, with possibly larger exponents.  They also retain a finite
bound on the cohomological amplitude of $C^*(\mathcal M_n;E)$.

For clarity, the primary part can be taken on spectra and not merely on
their homotopy groups.  A bounded Postnikov spectrum with bounded
$p$-primary exponent is killed, as an object, by a sufficiently large
power of $p$.  A polynomial in an endomorphism that vanishes on every
homotopy group becomes zero as a map after a bounded power, by the
Postnikov filtration.  Increase the two primary exponents accordingly.
The Bezout identity for the coprime factors $(T-1)^N$ and $Q(T)^N$
modulo that power of $p$ then gives complementary idempotents.  Splitting
them defines $C^*(\mathcal M_n;E)_{(1)}$ functorially.  Equivalently,
localization of the $\gamma$-action by $Q(\gamma)$ selects this summand.
The construction is exact on fibre sequences and is independent of the
larger exponents used to express the two primary factors.

The finite units $\eta_{X,G,E}$ are $\gamma$-equivariant by their
base-change naturality, and their continuous sources have trivial
$\gamma$-action.  They therefore land in this summand.  Applying
coefficient exact sequences to these same units extends
\eqref{eq:const-weight-zero} from $H\Fp$ to finite $p$-primary
Eilenberg--Mac Lane spectra.  Induction through the finite Postnikov
tower of $E$ proves the first two equivalences in
\eqref{eq:const-postnikov-transport}.  The last is
Lemma~\ref{lem:const-building}.  This proof uses the maps of cochain
spectra throughout, and their coefficient naturality gives Hurewicz
naturality.  Product compatibility means compatibility with a specified
coefficient pairing $E\wedge F\to E'$.  The termwise cup maps for
$\eta$ respect this pairing, and a tensor product of two unipotent
arithmetic actions is again unipotent, so they preserve the
generalized-one parts.

In particular, write the continuous sphere cochains as
\begin{equation}
 C^*_{\mathrm{cts}}(G;S)
   =\varprojlim_{r,N}
      C^*_{\mathrm{cts}}
          (G;\tau_{\leq N}(S/p^r)).
 \label{eq:const-sphere-cochains}
\end{equation}
Every displayed coefficient has finite $p$-primary homotopy groups.
Take the inverse limit of the already constructed finite units and
equivalences, defining the middle generalized-one object by this same
limit.  Every geometric pullback of a sphere class means this inverse
limit of the finite constant-section pullbacks.  We do not interchange
an eigenspace localization with an unrestricted inverse limit.  On the
classifying groups, finite cohomological
dimension controls the Postnikov limit in each degree; finite
coefficient cohomology and the finite building comparison control the
coefficient limit.  This agrees with continuous totalization of
termwise $p$-completed locally constant sphere cochains.  Consequently
a sphere-cochain representative on $K_n$ transports to one on $P_n$
with its specified mod-$p$ Hurewicz image.

\emph{The action on a fixed connected quotient.}
It remains to prove \eqref{eq:const-block-action}.  Fix a length-one
quotient $V_t\to i_*\mathcal L$ and a rational frame, with determinant
lattice, of its slope-zero kernel $W_t$.  On the full connected frame
space the universal sequence is
\[
 0\longrightarrow L_m\otimes\cO
      \longrightarrow V_n'\longrightarrow V_t\longrightarrow0.
\]
Modify it by the fixed quotient of $V_t$.  On $[Z/K_m]$ this produces
\begin{equation}
 \begin{gathered}
 0\longrightarrow W_n\longrightarrow V_n'
            \longrightarrow i_*\mathcal L\longrightarrow0,\\
 0\longrightarrow L_m\otimes\cO
            \longrightarrow W_n\longrightarrow W_t\longrightarrow0.
 \end{gathered}
 \label{eq:const-modified-extension}
\end{equation}
The first line defines a map $f:[Z/K_m]\to\mathcal M_n$ whose
positive-bundle classifying map is $b$.  All bundles in the second line
have slope zero.  Exactness of the slope-zero/local-system equivalence
therefore identifies its kernel with the actual rank-$m$ Tate local
system $L_m\otimes\Qp$, and its quotient with the fixed framed
$\Qp^t$.  Locally on the pro-\'{e}tale, hence on the $v$-site, this
sequence admits rational frames and splittings.

Its permitted frame changes form
\begin{equation}
 H_{m,t}=
 \left\{\begin{pmatrix}g&v\\0&1_t\end{pmatrix}:
          g\in K_m,\ v\in M_{m,t}(\Qp)\right\}\subset J_n.
 \label{eq:const-upper-block}
\end{equation}
The determinant condition is the product of the determinant lattice of
$L_m$ and the fixed one of $W_t$, using the matching convention in
Lemma~\ref{lem:const-modification-stack}.  The quotient
$H_{m,t}\to K_m$ has its displayed block-diagonal section.  The
composition $[Z/K_m]\to BH_{m,t}\to BK_m$ is exactly $q$.

The additive group $M_{m,t}(\Qp)$ has no positive continuous cohomology
with constant finite $p$-primary coefficients.  Indeed, exhaust it by
$p^{-r}\Zp^{mt}$.  The cohomology of these compact lattices with
$\Fp$-coefficients is the exterior algebra on their continuous duals.
Restriction to the preceding lattice multiplies each degree-one
generator by $p$, and is zero in every positive degree.  In degree
zero it is the identity.  The countable inverse-limit sequence has no
derived-limit error: these cohomology systems are Mittag--Leffler, and
restriction of continuous cochain terms is surjective because a
compact open subset admits extension of locally constant functions.
Thus only degree zero survives.  Coefficient exact sequences and
Postnikov induction give the same assertion for every finite
Postnikov $E$.

For the extension descent, the same lattice argument applies with an
extra compact profinite parameter $T$.  At $\Fp$ coefficients, finite
clopen rectangle refinements give
\[
 \operatorname{LC}(T\times(p^{-r}\Zp^{mt})^\bullet;\Fp)
  =\operatorname*{colim}_{T\twoheadrightarrow T_i}
       \prod_{t\in T_i}C^\bullet_{\mathrm{cts}}(p^{-r}\Zp^{mt};\Fp).
\]
Filtered colimits and finite products are exact, so its vertical
cohomology is the locally constant $T$-family of the compact-lattice
cohomology just computed.  The lattice restrictions are still the
identity in degree zero and zero in positive degrees, and the cochain
term restrictions are surjective by extension by zero from the clopen
subset $T\times(p^{-r}\Zp^{mt})^a$.  The same Mittag--Leffler argument
therefore applies.  The lattices are $K_m$-stable; take $T=K_m^b$ in
the outer bar for
$1\to M_{m,t}(\Qp)\to H_{m,t}\to K_m\to1$.  The resulting vertical
equivalences commute with the outer faces.  Totalizing this external
double bar, and then using coefficient and Postnikov induction, shows
that inflation from $K_m$ is an equivalence on the finite constant
cochains of $H_{m,t}$.  Its inverse is restriction along the displayed
section.

For a $J_n$-class $c_J$ whose restriction to $K_n$ is $c_K$, that
section followed by $H_{m,t}\to J_n$ is the upper-left block inclusion
through $K_n$.  Hence
\[
 \operatorname{res}^{J_n}_{H_{m,t}}c_J
   =\operatorname{inf}^{H_{m,t}}_{K_m}
       \operatorname{res}^{K_n}_{K_m}c_K.
\]
The $H_{m,t}$-torsor above induces the split-frame torsor of $f$ along
$H_{m,t}\to J_n$ and the torsor of $q$ along $H_{m,t}\to K_m$.
The finite torsor naturality of $\eta$ therefore identifies the pullback
of the left side with the $J_n$ structural image pulled along $f$, and
the pullback of the right side with the $K_m$ structural image pulled
along $q$.

The structural images of corresponding $c_P$ and $c_J$ in
$\mathcal M_n$ agree under
\eqref{eq:const-postnikov-transport} inside the generalized-one
summand itself.  Pull that equality along $f$.  Its positive-bundle
torsor is $b$ by \eqref{eq:const-modified-extension}, and the preceding
torsor calculation identifies its split-frame side.  Over the base
$\mathcal B$ of this diagram the resulting equality is
\[
 b^*\bigl(\eta_{\mathcal B,P_n,E}(c_P)\bigr)
   =q^*\bigl(\eta_{\mathcal B,K_m,E}
       (\operatorname{res}^{K_n}_{K_m}c_K)\bigr),
\]
which is \eqref{eq:const-block-action} under the classifying-pullback
convention.  The same equality persists under
\eqref{eq:const-sphere-cochains}.  For ordinary cohomology it is an
equality of the classes that act by cup product on the supported
complex on $Z$.  Once this action is identified with the $K_m$-action,
the $K_m$-equivariant map to the ambient affine support calculation
preserves it.  There is no need to extend the $BP_n$-classifying map
away from the independent locus.
\end{proof}
